\documentclass{resume} % 使用原 Overleaf 模板中的 resume.cls
\usepackage[UTF8]{ctex}
\usepackage[dvipsnames]{xcolor}
\usepackage{hyperref}
\usepackage[left=0.75in,top=0.55in,right=0.75in,bottom=0.35in]{geometry}

\name{冯恺睿}
\address{电话：+86 185 9965 1229 \quad 邮箱：\href{mailto:3405402279@qq.com}{3405402279@qq.com}}

\definecolor{CarnegieMellonRed}{RGB}{196,18,48}
\renewenvironment{rSection}[1]{
\sectionskip
\textcolor{CarnegieMellonRed}{\MakeUppercase{#1}}
\sectionlineskip
\hrule
\begin{list}{}{\setlength{\leftmargin}{1.5em}}
\item[]
}{\end{list}}

\begin{document}

\begin{rSection}{教育背景}
\textbf{浙江大学} \hfill \textit{预计 2027 年 6 月毕业}\\
计算机科学与技术学院，信息安全专业，工学学士\\
GPA：3.96/4.30 \quad 加权平均分：87.33/100 \quad 雅思：6.5
\end{rSection}

\begin{rSection}{安全相关核心课程}
\textbf{安全相关课程加权平均分：92.1/100（18.5 学分）}\\
\small 软件安全原理（99）；软件安全理论与实践（92）；网络安全理论与实践（92）；人工智能安全（89）；人工智能安全与伦理（89）；高级密码学（92）；系统安全：原理与实践（97）；无线与物联网安全（84）；隐私、安全与数字货币（95）。\normalsize
\end{rSection}

\begin{rSection}{系统与安全项目经历}
\textbf{Linux 内核安全分析与利用} \hfill \textit{独立安全实践}\\
\begin{itemize}
    \item 通过源码阅读分析内核染色技术。
\end{itemize}

\textbf{gdb-helper} \hfill \textit{开源调试工具}\\
\begin{itemize}
    \item 在已有开源 GDB 插件中接入 DeepSeek API，将命令行中的自然语言调试需求转化为多步原生 GDB 调试流程。\url{https://github.com/fkrcarry/gdb-helper}
\end{itemize}

\textbf{运行于自研 CPU 的 RISC-V 操作系统} \hfill \textit{浙江大学计算机系统贯通课程}\\
\begin{itemize}
    \item 独立实现运行于自研 CPU 的 RISC-V 操作系统；\textbf{SYS1：}使用 HDL 设计、实现并验证单周期 RISC-V CPU。
    \item \textbf{SYS2：}实现支持中断和进程上下文切换的简易内核；实现支持五级流水线、分支预测、总线接口、转发、一级缓存和特权指令的 RV64I CPU。
    \item \textbf{SYS3（Lab 1--5 及 xpart MMU）：}实现启动和时钟中断、抢占式内核线程、虚拟内存、用户态和系统调用、缺页、进程创建及 MMU 功能。
    \item 获常瑞老师、申文博老师颁发的\textit{优秀作品奖}；约 150 名信息安全和图灵班学生中仅 9 人获奖。
\end{itemize}

\textbf{C++ RISC-V 编译器} \hfill \textit{独立系统项目}\\
\begin{itemize}
    \item 使用 C++ 开发面向 RISC-V 的编译器。
\end{itemize}

\textbf{基于移动端 GUI 智能体的 iOS 变身 App 检测系统} \hfill \textit{许海涛老师指导}\\
\begin{itemize}
    \item 开发基于移动端 GUI 智能体的 iOS 应用检测系统，识别交互中改变行为或呈现方式的应用。
\end{itemize}

\textbf{MCP Server 安全漏洞挖掘} \hfill \textit{独立安全研究}\\
\begin{itemize}
    \item 发现并公开披露 \texttt{imprvhub/}\allowbreak\texttt{mcp-browser-agent}（至 0.8.0 版本）的 SSRF 漏洞，根因是 \texttt{CallToolRequestSchema} 中用户可控字段；CVE-2026-5607：\url{https://vuldb.com/source_cve/355398}
\end{itemize}

\textbf{其他安全研究} \hfill \textit{浙江大学}\\
\begin{itemize}
    \item 在任奎老师、秦湛老师指导下参与代码大模型红队研究（相关专利申请中）；在杜文亮老师指导下参与 SEED Labs 安全实验。
\end{itemize}
\end{rSection}

\begin{rSection}{助教经历}
\textbf{浙江大学助教：}计算机系统 I（SYS1）、计算机系统 II（SYS2）、软件安全原理与实践。\\
负责理论与实验讲解、实验发布和验收、作业及 Quiz 批改。
\end{rSection}

\begin{rSection}{荣誉与竞赛}
\begin{itemize}
    \item 浙江大学校级标兵（两次）
    \item 全国大学生信息安全竞赛二等奖
    \item 浙江省大学生信息安全竞赛一等奖
    \item 2025 ZJUCTF Pwn 方向校内第三名
    \item 浙江大学 CTF 网站 Pwn 方向校内第 4 名
    \item 计算机系统贯通课程优秀作品奖
\end{itemize}
\end{rSection}

\begin{rSection}{专业技能}
\begin{tabular}{@{} >{\bfseries}l @{\hspace{1.5ex}} p{0.70\textwidth}}
编程语言 & C、C++、Python、SystemVerilog \\
系统与工具 & Linux 命令行、Git、Docker、\LaTeX、QEMU 内核调试、二进制程序调试、内核漏洞复现与分析、Fuzzing、Linux 内核模块/驱动、内核提权利用 \\
\end{tabular}
\end{rSection}

\end{document}
